% Options for packages loaded elsewhere
\PassOptionsToPackage{unicode}{hyperref}
\PassOptionsToPackage{hyphens}{url}
\PassOptionsToPackage{dvipsnames,svgnames,x11names}{xcolor}
\documentclass[
  10pt,
  a4paper,
]{article}
\usepackage{xcolor}
\usepackage[margin=22mm]{geometry}
\usepackage{amsmath,amssymb}
\setcounter{secnumdepth}{-\maxdimen} % remove section numbering
\usepackage{iftex}
\ifPDFTeX
  \usepackage[T1]{fontenc}
  \usepackage[utf8]{inputenc}
  \usepackage{textcomp} % provide euro and other symbols
\else % if luatex or xetex
  \usepackage{unicode-math} % this also loads fontspec
  \defaultfontfeatures{Scale=MatchLowercase}
  \defaultfontfeatures[\rmfamily]{Ligatures=TeX,Scale=1}
\fi
\usepackage{lmodern}
\ifPDFTeX\else
  % xetex/luatex font selection
\fi
% Use upquote if available, for straight quotes in verbatim environments
\IfFileExists{upquote.sty}{\usepackage{upquote}}{}
\IfFileExists{microtype.sty}{% use microtype if available
  \usepackage[]{microtype}
  \UseMicrotypeSet[protrusion]{basicmath} % disable protrusion for tt fonts
}{}
\makeatletter
\@ifundefined{KOMAClassName}{% if non-KOMA class
  \IfFileExists{parskip.sty}{%
    \usepackage{parskip}
  }{% else
    \setlength{\parindent}{0pt}
    \setlength{\parskip}{6pt plus 2pt minus 1pt}}
}{% if KOMA class
  \KOMAoptions{parskip=half}}
\makeatother
\usepackage{color}
\usepackage{fancyvrb}
\newcommand{\VerbBar}{|}
\newcommand{\VERB}{\Verb[commandchars=\\\{\}]}
\DefineVerbatimEnvironment{Highlighting}{Verbatim}{commandchars=\\\{\}}
% Add ',fontsize=\small' for more characters per line
\newenvironment{Shaded}{}{}
\newcommand{\AlertTok}[1]{\textcolor[rgb]{1.00,0.00,0.00}{\textbf{#1}}}
\newcommand{\AnnotationTok}[1]{\textcolor[rgb]{0.38,0.63,0.69}{\textbf{\textit{#1}}}}
\newcommand{\AttributeTok}[1]{\textcolor[rgb]{0.49,0.56,0.16}{#1}}
\newcommand{\BaseNTok}[1]{\textcolor[rgb]{0.25,0.63,0.44}{#1}}
\newcommand{\BuiltInTok}[1]{\textcolor[rgb]{0.00,0.50,0.00}{#1}}
\newcommand{\CharTok}[1]{\textcolor[rgb]{0.25,0.44,0.63}{#1}}
\newcommand{\CommentTok}[1]{\textcolor[rgb]{0.38,0.63,0.69}{\textit{#1}}}
\newcommand{\CommentVarTok}[1]{\textcolor[rgb]{0.38,0.63,0.69}{\textbf{\textit{#1}}}}
\newcommand{\ConstantTok}[1]{\textcolor[rgb]{0.53,0.00,0.00}{#1}}
\newcommand{\ControlFlowTok}[1]{\textcolor[rgb]{0.00,0.44,0.13}{\textbf{#1}}}
\newcommand{\DataTypeTok}[1]{\textcolor[rgb]{0.56,0.13,0.00}{#1}}
\newcommand{\DecValTok}[1]{\textcolor[rgb]{0.25,0.63,0.44}{#1}}
\newcommand{\DocumentationTok}[1]{\textcolor[rgb]{0.73,0.13,0.13}{\textit{#1}}}
\newcommand{\ErrorTok}[1]{\textcolor[rgb]{1.00,0.00,0.00}{\textbf{#1}}}
\newcommand{\ExtensionTok}[1]{#1}
\newcommand{\FloatTok}[1]{\textcolor[rgb]{0.25,0.63,0.44}{#1}}
\newcommand{\FunctionTok}[1]{\textcolor[rgb]{0.02,0.16,0.49}{#1}}
\newcommand{\ImportTok}[1]{\textcolor[rgb]{0.00,0.50,0.00}{\textbf{#1}}}
\newcommand{\InformationTok}[1]{\textcolor[rgb]{0.38,0.63,0.69}{\textbf{\textit{#1}}}}
\newcommand{\KeywordTok}[1]{\textcolor[rgb]{0.00,0.44,0.13}{\textbf{#1}}}
\newcommand{\NormalTok}[1]{#1}
\newcommand{\OperatorTok}[1]{\textcolor[rgb]{0.40,0.40,0.40}{#1}}
\newcommand{\OtherTok}[1]{\textcolor[rgb]{0.00,0.44,0.13}{#1}}
\newcommand{\PreprocessorTok}[1]{\textcolor[rgb]{0.74,0.48,0.00}{#1}}
\newcommand{\RegionMarkerTok}[1]{#1}
\newcommand{\SpecialCharTok}[1]{\textcolor[rgb]{0.25,0.44,0.63}{#1}}
\newcommand{\SpecialStringTok}[1]{\textcolor[rgb]{0.73,0.40,0.53}{#1}}
\newcommand{\StringTok}[1]{\textcolor[rgb]{0.25,0.44,0.63}{#1}}
\newcommand{\VariableTok}[1]{\textcolor[rgb]{0.10,0.09,0.49}{#1}}
\newcommand{\VerbatimStringTok}[1]{\textcolor[rgb]{0.25,0.44,0.63}{#1}}
\newcommand{\WarningTok}[1]{\textcolor[rgb]{0.38,0.63,0.69}{\textbf{\textit{#1}}}}
\usepackage{longtable,booktabs,array}
\newcounter{none} % for unnumbered tables
\usepackage{calc} % for calculating minipage widths
% Correct order of tables after \paragraph or \subparagraph
\usepackage{etoolbox}
\makeatletter
\patchcmd\longtable{\par}{\if@noskipsec\mbox{}\fi\par}{}{}
\makeatother
% Allow footnotes in longtable head/foot
\IfFileExists{footnotehyper.sty}{\usepackage{footnotehyper}}{\usepackage{footnote}}
\makesavenoteenv{longtable}
\usepackage{graphicx}
\makeatletter
\newsavebox\pandoc@box
\newcommand*\pandocbounded[1]{% scales image to fit in text height/width
  \sbox\pandoc@box{#1}%
  \Gscale@div\@tempa{\textheight}{\dimexpr\ht\pandoc@box+\dp\pandoc@box\relax}%
  \Gscale@div\@tempb{\linewidth}{\wd\pandoc@box}%
  \ifdim\@tempb\p@<\@tempa\p@\let\@tempa\@tempb\fi% select the smaller of both
  \ifdim\@tempa\p@<\p@\scalebox{\@tempa}{\usebox\pandoc@box}%
  \else\usebox{\pandoc@box}%
  \fi%
}
% Set default figure placement to htbp
\def\fps@figure{htbp}
\makeatother
\setlength{\emergencystretch}{3em} % prevent overfull lines
\providecommand{\tightlist}{%
  \setlength{\itemsep}{0pt}\setlength{\parskip}{0pt}}
\usepackage{bookmark}
\IfFileExists{xurl.sty}{\usepackage{xurl}}{} % add URL line breaks if available
\urlstyle{same}
\hypersetup{
  colorlinks=true,
  linkcolor={blue},
  filecolor={Maroon},
  citecolor={Blue},
  urlcolor={blue},
  pdfcreator={LaTeX via pandoc}}

\author{}
\date{}

\begin{document}

\section{Discovery of Finite-Order Resonance in Iterated Exchange
Scattering}\label{discovery-of-finite-order-resonance-in-iterated-exchange-scattering}

\subsection{Identifying Sharp Peaks near Fine-Structure-Constant Inverse
Values 137 and 128 with a Reproducible Wave-Packet
Model}\label{identifying-sharp-peaks-near-fine-structure-constant-inverse-values-137-and-128-with-a-reproducible-wave-packet-model}

\textbf{Version:} English complete manuscript v1.0\\
\textbf{Date:} July 18, 2026\\
\textbf{Author:} Noriaki Kihara\\
\textbf{Series position:} Wave Information Readout series; paper
identifying the origin of exchange-weight peaks\\
\textbf{Version DOI:} 10.5281/zenodo.21421367\\
\textbf{Concept DOI:} 10.5281/zenodo.21421366

\begin{center}\rule{0.5\linewidth}{0.5pt}\end{center}

\subsection{Abstract}\label{abstract}

This study began with the explicit aim of testing whether an
exchange-scattering model of a closed phase system could explain values
associated with the low- and high-energy fine-structure constant. In the
preceding study, the exchange-scattering coefficient \(R\) was swept in
System A, which reads localization transfer between wave packets, and
System B, which reads a white-cat/black-cat/gray-cat metastable
interface. Sharp peaks were observed near the regions corresponding to
\(\alpha^{-1}\approx137\) and \(\alpha^{-1}\approx128\), but the reason
why peaks formed at those positions remained unresolved.

This paper reports a reanalysis of the full-range sweep, root-centered
precision sweeps, operator-eigenvalue analysis, and 50- and 80-digit
calculations performed to identify the cause. The analysis shows that
the observed peaks do not arise from the assumed
electromagnetic-coupling readout alone. They are generated by exact
recurrence of the two-channel exchange-scattering operator

\[
U_R=
\begin{pmatrix}
r&t\\
t&r
\end{pmatrix},
\]

whose antisymmetric eigenvalue returns to unity after a finite number of
iterations:

\[
U_R^n=I.
\]

The finite-order roots are

\[
\boxed{
R_{n,m}=\cos^2\left(\frac{\pi m}{n}\right)
}.
\]

The principal low-energy peak is \(R_{124,23}\), and the former
high-energy principal candidate is \(R_{122,23}\). A separate even-order
root \(R_{620,117}\) lies near the high-energy value
\(\alpha^{-1}(M_Z^2)=128.946\pm0.015\) and gives

\[
N(R_{620,117})
=\frac{4\pi}{(1-R_{620,117})^2}
=128.947864735670559.
\]

Its difference from the quoted value is \(0.0018647\), or
\(0.1243\sigma\).

The 50-digit calculations produced depths near 52, while the 80-digit
calculations produced depths from 82 to 83. At an exact finite-order
root, the residual vanishes in ideal arithmetic, so the depth is
formally infinite. This is not a divergence of energy or amplitude; it
means that the resonance line becomes infinitely sharp.

The principal result is not only the analytic root formula. Two
numerical models originally constructed for different
purposes---wave-packet localization transfer and observation selection
at a metastable two-state interface---share the same complex
exchange-scattering kernel and reproduce operator-intrinsic finite-order
resonance. Exact recurrence can therefore be tested as an iterated
wave-packet interaction calculation rather than remaining only an
abstract operator statement.

After this analysis, a literature search revealed connections to exact
state-independent Floquet recurrence from cyclotomic spectra, Farey
organization of resonance diagrams, the decomposition of \(E_8\) into 31
orthogonal eight-dimensional subspaces, and studies relating exceptional
algebras to the fine-structure constant. These works did not motivate
the numerical experiment; they were identified only after the
finite-order roots had been found.

We initially expected this numerical experiment to explain \(\alpha\).
What we directly found was a different finite-order resonance phenomenon
connected to a generally known class of exact recurrences. It remains
unresolved why such sharp finite-order roots occur near the two physical
values of \(\alpha^{-1}\). Determining whether this proximity is an
arithmetic coincidence or follows from four-dimensional discrete
geometry or another structural necessity is the central task of the next
stage.

\textbf{Keywords:} fine-structure constant, exchange scattering,
finite-order resonance, quantum recurrence, Floquet operator,
wave-packet localization, metastable state, \(E_8\), reproducible
computation

\begin{center}\rule{0.5\linewidth}{0.5pt}\end{center}

\section{1. Research Motivation}\label{research-motivation}

\subsection{1.1 A Search Deliberately Aimed at the Fine-Structure
Constant}\label{a-search-deliberately-aimed-at-the-fine-structure-constant}

The motivation of this study must not be obscured. We explicitly
expected that the exchange-scattering structure of a closed phase system
might explain a value corresponding to the fine-structure constant.

We define the exchange weight by

\[
G_R:=1-R
\]

and assume the correspondence

\[
G_R=e
\]

with the electromagnetic coupling amplitude in the standard theory,

\[
e=\sqrt{4\pi\alpha}.
\]

This is the correspondence hypothesis of the present study. We further
define the readout

\[
N(R):=\alpha^{-1}.
\]

Since

\[
(1-R)^2=4\pi\alpha,
\]

we obtain

\[
\boxed{
N(R)=\frac{4\pi}{(1-R)^2}
}.
\]

Under this correspondence hypothesis, exchange-scattering coefficients
can be assigned to the low-energy value

\[
\alpha^{-1}(0)=137.035999177(21)
\]

and the high-energy value

\[
\alpha^{-1}(M_Z^2)=128.946\pm0.015.
\]

Thus, the region \(R\approx0.7\) was not selected accidentally. We
examined whether the effective-response region found in the
localization-transfer and metastable-interface experiments could connect
to the exchange weight corresponding to the electromagnetic coupling
amplitude.

\subsection{1.2 System A: Wave-Packet Localization
Transfer}\label{system-a-wave-packet-localization-transfer}

The first starting point is a localization-exchange model in which an
exchange-interference scattering matrix acts on wave-packet vectors
{[}1,2{]}. A is a low-localization wave and B is a localized wave
containing additional harmonics. We iterate

\[
\begin{pmatrix}
A'\\B'
\end{pmatrix}
=
U_R
\begin{pmatrix}
A\\B
\end{pmatrix}.
\]

The readouts are

\[
L,\qquad N_{\mathrm{eff}},\qquad B_{\mathrm{to}A},
\]

where \(L\) is localization, \(N_{\mathrm{eff}}\) is the effective
harmonic order, and \(B_{\mathrm{to}A}\) measures the transfer of the
initial B-side harmonic distribution into A.

At the endpoints of complete transmission \(R=0\) and complete
reflection \(R=1\), the distinction between the low- and high-order
waves was preserved. Under intermediate scattering, localization and
harmonic structure were periodically transferred between the two
channels. In particular, the localization and effective harmonic orders
of the two channels approached one another near \(R\approx0.70\).

This model was not originally designed to search for exact recurrence.
It was constructed to numerically investigate an interaction in which a
wave packet appears to become localized through exchange with another
wave packet.

\subsection{1.3 System B: White-Cat, Black-Cat, and Gray-Cat Metastable
Interface}\label{system-b-white-cat-black-cat-and-gray-cat-metastable-interface}

The second starting point is the white-cat/black-cat/gray-cat model,
which reads the distribution and observation selection of two states A
and B {[}3{]}. From complex amplitudes \(a\) and \(b\), define

\[
p_A=|a|^2,\qquad
p_B=|b|^2,\qquad
S=p_A-p_B.
\]

\begin{itemize}
\tightlist
\item
  \(S\approx+1\): white-cat state A.
\item
  \(S\approx-1\): black-cat state B.
\item
  \(S\approx0\): gray state.
\end{itemize}

Among gray states, a state that fluctuates with small amplitude,
survives a weak C readout, and is selected into A or B by a strong D
observation is called a gray metastable phase.

This model was likewise not designed from the outset to reproduce the
fine-structure constant or Floquet recurrence. It was constructed to
distinguish a readout that does not destroy a metastable mixed state
from strong observational selection.

\subsection{1.4 The Question Left by the Previous
Paper}\label{the-question-left-by-the-previous-paper}

The paper of July 15, 2026 {[}4{]} connected Systems A and B through
their common exchange-scattering coefficient \(R\). It found that:

\begin{itemize}
\tightlist
\item
  with an additional harmonic, the response concentrates in the region
  corresponding to 137;
\item
  the candidate position barely moves when the harmonic type, amplitude,
  phase, or wavelength is changed; and
\item
  the full System B sweep contains deep peaks near the regions read as
  137 and 128.
\end{itemize}

The following question, however, remained unresolved:

\begin{quote}
Why do sharp peaks form at those positions?
\end{quote}

This paper reports the investigation undertaken to answer that question
and the resulting identification of the cause.

\begin{center}\rule{0.5\linewidth}{0.5pt}\end{center}

\section{2. Computational Model and
Readout}\label{computational-model-and-readout}

\subsection{2.1 Exchange-Scattering
Coefficients}\label{exchange-scattering-coefficients}

For a reflection probability \(R\in[0,1]\), define

\[
\theta(R):=\arcsin\sqrt R.
\]

The transmission and reflection amplitudes are

\[
t=e^{i\theta}\cos\theta,
\qquad
r=-ie^{i\theta}\sin\theta.
\]

Therefore,

\[
|t|^2=1-R,
\qquad
|r|^2=R.
\]

\subsection{2.2 Two-Channel Exchange
Operator}\label{two-channel-exchange-operator}

One scattering iteration is defined by

\[
\boxed{
\begin{pmatrix}
A_{j+1}\\
B_{j+1}
\end{pmatrix}
=
U_R
\begin{pmatrix}
A_j\\
B_j
\end{pmatrix},
\qquad
U_R=
\begin{pmatrix}
r&t\\
t&r
\end{pmatrix}
}.
\]

Because the coefficients satisfy

\[
|r|^2+|t|^2=1,
\qquad
r^*t+t^*r=0,
\]

we have

\[
U_R^\dagger U_R=I.
\]

Thus, the exact fixed-\(R\) map is unitary and does not dissipatively
attract the state norm to a point.

The ``concentration'\,' at a peak in this study is not dissipative
contraction in state space. It is resonance selection in \(R\)-parameter
space, where a closed-orbit condition concentrates at particular values.
The increase of wave-packet localization observed in System A is a
redistribution of localization by unitary interference, not contraction
of the total norm.

\subsection{2.3 Initial State of System
B}\label{initial-state-of-system-b}

Let

\[
A_0=\sqrt{\frac12+s_0},
\qquad
B_0=\sqrt{\frac12-s_0}\,e^{i\phi}.
\]

The principal sweep uses

\[
s_0=0.01,\qquad \phi=0.
\]

The initial population difference is therefore

\[
S_0=|A_0|^2-|B_0|^2=0.02.
\]

\subsection{2.4 Gray Error and Depth}\label{gray-error-and-depth}

At each time, define the population difference by

\[
S_j=
\frac{|A_j|^2-|B_j|^2}
{|A_j|^2+|B_j|^2}.
\]

For a prefix of length \(k\), define

\[
\overline S_k
=\frac1k\sum_{j=0}^{k-1}S_j,
\]

\[
A_k
=\frac{\max_{j<k}S_j-\min_{j<k}S_j}{2},
\]

and

\[
D_k
=\left|
\overline S_{k,\mathrm{second}}
-\overline S_{k,\mathrm{first}}
\right|.
\]

The gray error is

\[
\varepsilon_k
=|\overline S_k|
+|A_k-0.02|
+D_k,
\]

and the depth is

\[
\boxed{
d_k=-\log_{10}\varepsilon_k
}.
\]

A periodic orbit becomes a deeper peak when its mean approaches zero,
its amplitude approaches \(0.02\), and the drift between the first and
second halves approaches zero.

\begin{center}\rule{0.5\linewidth}{0.5pt}\end{center}

\section{3. Full-Range Sweep in the Previous
Study}\label{full-range-sweep-in-the-previous-study}

\subsection{3.1 Sweep Conditions}\label{sweep-conditions}

Using the minimal System B implementation v5, we uniformly swept 158,625
points over

\[
R_{\min}=0.68660290255614798,
\]

\[
R_{\max}=0.70246536756843059,
\]

with

\[
\Delta R=10^{-7}.
\]

The same scattering update, gray error, and stopping condition were
applied to every \(R\). The sweep evaluation formula contains neither
\(\alpha\), 137, 128, nor a finite-order root.

\subsection{3.2 Full-Range Sweep Figure}\label{full-range-sweep-figure}

\pandocbounded{\includegraphics[keepaspectratio]{../20260715/system_B_full_R_sweep_full_range_depth_v1.png}}

\textbf{Figure 1.} Full System B sweep obtained on July 15, 2026. The
horizontal axis is \(R\), and the vertical axis is depth normalized
within the candidate bands. The deepest candidate 1 appeared at
\(R=0.697177902556148\), and candidate 2 appeared at
\(R=0.688363902556148\).

The full-range figure showed that the response was not a single peak but
a collection of narrow candidate bands. At that time, candidate 1 was
read as a peak corresponding to the low-energy value
\(\alpha^{-1}(0)\approx137\), while candidate 2 was read as a peak near
the high-energy value \(\alpha^{-1}(M_Z^2)\approx128\).

However, a spacing of \(10^{-7}\) could not determine the true center,
width, limiting depth, or generating mechanism of each peak.

\subsection{3.3 Original Candidate
Values}\label{original-candidate-values}

{\def\LTcaptype{none} % do not increment counter
\begin{longtable}[]{@{}
  >{\raggedright\arraybackslash}p{(\linewidth - 6\tabcolsep) * \real{0.2143}}
  >{\raggedleft\arraybackslash}p{(\linewidth - 6\tabcolsep) * \real{0.2857}}
  >{\raggedleft\arraybackslash}p{(\linewidth - 6\tabcolsep) * \real{0.2857}}
  >{\raggedright\arraybackslash}p{(\linewidth - 6\tabcolsep) * \real{0.2143}}@{}}
\toprule\noalign{}
\begin{minipage}[b]{\linewidth}\raggedright
Candidate
\end{minipage} & \begin{minipage}[b]{\linewidth}\raggedleft
\(R\) in the full sweep
\end{minipage} & \begin{minipage}[b]{\linewidth}\raggedleft
\(N(R)\)
\end{minipage} & \begin{minipage}[b]{\linewidth}\raggedright
Interpretation at the time
\end{minipage} \\
\midrule\noalign{}
\endhead
\bottomrule\noalign{}
\endlastfoot
1 & 0.697177902556148 & 137.036020287643 & Near low-energy 137 \\
2 & 0.688363902556148 & 129.394062925467 & Candidate near high-energy
128 \\
\end{longtable}
}

Candidate 2 was not itself \(128.946\). Nevertheless, because it was the
second principal peak in the same full-range figure, its origin required
identification.

\begin{center}\rule{0.5\linewidth}{0.5pt}\end{center}

\section{4. Root-Centered Precision
Sweeps}\label{root-centered-precision-sweeps}

\subsection{\texorpdfstring{4.1 Low-Energy Side:
\(R_{124,23}\)}{4.1 Low-Energy Side: R\_\{124,23\}}}\label{low-energy-side-r_12423}

Candidate 1 from the full sweep was magnified while reducing the spacing
successively to

\[
10^{-8},\quad10^{-10},\quad10^{-12}.
\]

The peak center converged to

\[
\boxed{
R_{124,23}
=\cos^2\left(\frac{23\pi}{124}\right)
=0.697177927556659290\ldots
}.
\]

\pandocbounded{\includegraphics[keepaspectratio]{root_centered_delta1e-12_v1/root_centered_delta1e-12_v1.png}}

\textbf{Figure 2.} A \(\Delta R=10^{-12}\) sweep with \(R_{124,23}\)
placed on the central grid point. The deepest point, analytic root, and
central grid point coincide. In double precision, the depth is 14.6586
and the error is \(2.19\times10^{-15}\).

Reading this root through the correspondence hypothesis
\(N(R)=4\pi/(1-R)^2\) gives

\[
N(R_{124,23})
=137.036042914598545\ldots.
\]

The CODATA 2022 value is

\[
\alpha^{-1}(0)=137.035999177(21),
\]

so the difference is

\[
4.37376\times10^{-5}.
\]

The two values are close but not identical.

\subsection{\texorpdfstring{4.2 Former High-Energy Candidate:
\(R_{122,23}\)}{4.2 Former High-Energy Candidate: R\_\{122,23\}}}\label{former-high-energy-candidate-r_12223}

Candidate 2 was magnified in the same format. Its center is

\[
\boxed{
R_{122,23}
=\cos^2\left(\frac{23\pi}{122}\right)
=0.688363946817592550\ldots
}.
\]

\pandocbounded{\includegraphics[keepaspectratio]{R122_23_root_centered_resolution_comparison_v1/R122_23_root_centered_resolution_comparison_v1.png}}

\textbf{Figure 3.} Comparison of \(\Delta R=10^{-10}\) and \(10^{-12}\)
sweeps centered on \(R_{122,23}\). At both resolutions, the analytic
root coincides with the deepest point. The depth is 14.8777 and the
error is \(1.33\times10^{-15}\).

The corresponding readout is

\[
N(R_{122,23})
=129.394099680960693\ldots.
\]

Thus, the direct identity of the second principal peak previously called
a candidate near high-energy 128 is the finite-order resonance
\(R_{122,23}\).

\subsection{\texorpdfstring{4.3 Near the High-Energy Observed Value:
\(R_{620,117}\)}{4.3 Near the High-Energy Observed Value: R\_\{620,117\}}}\label{near-the-high-energy-observed-value-r_620117}

Higher-order finite-order roots were examined near the \(R\) value
corresponding to

\[
\alpha^{-1}(M_Z^2)=128.946\pm0.015.
\]

This search identified the even-order root

\[
\boxed{
R_{620,117}
=\cos^2\left(\frac{117\pi}{620}\right)
=0.687825191114145187\ldots
}.
\]

\pandocbounded{\includegraphics[keepaspectratio]{high_order_roots_centered_delta1e-12_v1/R620_117_root_centered_delta1e-12_v1.png}}

\textbf{Figure 4.} A \(\Delta R=10^{-12}\) sweep with \(R_{620,117}\) on
the central grid point. The deepest point is the analytic root itself.
The depth is 14.9693, the error is \(1.07\times10^{-15}\), and the best
prefix ends at \(2n-1=1239\).

This root gives

\[
N(R_{620,117})
=128.947864735670559\ldots.
\]

The difference from the observed value is

\[
128.947864735670559-128.946
=0.001864735670559,
\]

and hence

\[
\frac{0.001864735670559}{0.015}
=0.1243157,
\]

or \(0.1243\sigma\).

The previous full-range sweep used at most 1024 steps. The gray metric
for \(R_{620,117}\) requires

\[
2n=1240
\]

samples to read two complete periods. Therefore, this root could not
appear as a complete peak within the earlier time window. The root near
the high-energy observed value was found later not only because of
spatial resolution along \(R\), but also because the iteration length in
the time direction had been insufficient.

\subsection{\texorpdfstring{4.4 Odd-Order Control:
\(R_{567,107}\)}{4.4 Odd-Order Control: R\_\{567,107\}}}\label{odd-order-control-r_567107}

In the same observed region,

\[
R_{567,107}
=0.687816908934256712\ldots
\]

gives

\[
N(R_{567,107})
=128.941022883719000\ldots.
\]

However, \(n=567\) is odd, so the discrete orbit does not contain the
opposite gray-metric extremum \(-0.02\) as a grid point. Thus, even
though the operator satisfies

\[
U_R^{567}=I,
\]

the amplitude readout retains the deficit

\[
0.01\left(1-\cos\frac{\pi}{567}\right)
=1.5349792919\times10^{-7}.
\]

The measured deficit is

\[
1.5349792966\times10^{-7},
\]

in agreement with the prediction.

This control shows that finite-order operator recurrence and an
infinite-depth peak in the selected observable are not synonymous.
Operator recurrence holds regardless of the parity of \(n\), while the
present gray-amplitude readout produces a machine-precision peak at even
fundamental orders.

\begin{center}\rule{0.5\linewidth}{0.5pt}\end{center}

\section{5. Multiprecision Calculation and Infinitely Sharp
Peaks}\label{multiprecision-calculation-and-infinitely-sharp-peaks}

\subsection{5.1 50- and 80-Digit
Calculations}\label{and-80-digit-calculations}

We placed \(R_{124,23}\) and \(R_{620,117}\) on the central grid point
and calculated them with 50- and 80-digit arithmetic.

{\def\LTcaptype{none} % do not increment counter
\begin{longtable}[]{@{}
  >{\raggedright\arraybackslash}p{(\linewidth - 10\tabcolsep) * \real{0.1304}}
  >{\raggedleft\arraybackslash}p{(\linewidth - 10\tabcolsep) * \real{0.1739}}
  >{\raggedleft\arraybackslash}p{(\linewidth - 10\tabcolsep) * \real{0.1739}}
  >{\raggedleft\arraybackslash}p{(\linewidth - 10\tabcolsep) * \real{0.1739}}
  >{\raggedleft\arraybackslash}p{(\linewidth - 10\tabcolsep) * \real{0.1739}}
  >{\raggedleft\arraybackslash}p{(\linewidth - 10\tabcolsep) * \real{0.1739}}@{}}
\toprule\noalign{}
\begin{minipage}[b]{\linewidth}\raggedright
Root
\end{minipage} & \begin{minipage}[b]{\linewidth}\raggedleft
Precision
\end{minipage} & \begin{minipage}[b]{\linewidth}\raggedleft
\(\Delta R\)
\end{minipage} & \begin{minipage}[b]{\linewidth}\raggedleft
Central depth
\end{minipage} & \begin{minipage}[b]{\linewidth}\raggedleft
Central error
\end{minipage} & \begin{minipage}[b]{\linewidth}\raggedleft
Best step
\end{minipage} \\
\midrule\noalign{}
\endhead
\bottomrule\noalign{}
\endlastfoot
\(R_{124,23}\) & 50 digits & \(10^{-14}\) & 52.7036 &
\(1.98\times10^{-53}\) & 247 \\
\(R_{620,117}\) & 50 digits & \(10^{-14}\) & 52.6831 &
\(2.07\times10^{-53}\) & 1239 \\
\(R_{124,23}\) & 80 digits & \(10^{-16}\) & 82.1106 &
\(7.75\times10^{-83}\) & 247 \\
\(R_{620,117}\) & 80 digits & \(10^{-16}\) & 83.0581 &
\(8.75\times10^{-84}\) & 1239 \\
\end{longtable}
}

\pandocbounded{\includegraphics[keepaspectratio]{two_physical_roots_multiprecision_v1/80digit_delta1e-16_comparison_v1.png}}

\textbf{Figure 5.} An 80-digit multiprecision sweep of \(R_{124,23}\)
and \(R_{620,117}\). For both roots, the central grid point is the
deepest point, and the depth exceeds 80 as the arithmetic precision is
increased.

\subsection{5.2 The Theoretical Limit}\label{the-theoretical-limit}

At a root whose relatively prime representation \((m,n)\) has even
fundamental order \(n\), reading two complete periods \(2n\) gives
exactly

\[
\overline S=0,
\qquad
A=0.02,
\qquad
D=0.
\]

Therefore,

\[
\varepsilon=0,
\]

and

\[
d=-\log_{10}\varepsilon
\longrightarrow+\infty.
\]

The diverging quantity is gray depth, not wave energy or amplitude. The
height at the peak center increases with precision while the width
narrows. A finite-precision measurement or coarse \(R\) sweep can
therefore miss the peak easily.

\begin{center}\rule{0.5\linewidth}{0.5pt}\end{center}

\section{6. Mathematical Origin of the
Peaks}\label{mathematical-origin-of-the-peaks}

\subsection{6.1 Symmetric and Antisymmetric
Decomposition}\label{symmetric-and-antisymmetric-decomposition}

Decompose the two-channel state as

\[
X_j:=\frac{A_j+B_j}{\sqrt2},
\qquad
Y_j:=\frac{A_j-B_j}{\sqrt2}.
\]

The eigenvalues of the exchange operator are

\[
\lambda_s=r+t=1,
\]

and

\[
\lambda_a=r-t=-e^{2i\theta}.
\]

Therefore,

\[
X_j=X_0,
\qquad
Y_j=\lambda_a^jY_0.
\]

With the projection operators

\[
P_s=\frac12
\begin{pmatrix}
1&1\\1&1
\end{pmatrix},
\qquad
P_a=\frac12
\begin{pmatrix}
1&-1\\-1&1
\end{pmatrix},
\]

we obtain

\[
\boxed{
U_R^j=P_s+\lambda_a^jP_a
}.
\]

\subsection{6.2 Derivation of the Finite-Order
Roots}\label{derivation-of-the-finite-order-roots}

Let

\[
R_{n,m}
=\cos^2\left(\frac{\pi m}{n}\right),
\qquad
1\le m<\frac n2,
\qquad
\gcd(m,n)=1.
\]

Then

\[
\theta
=\arcsin\sqrt{R_{n,m}}
=\frac{\pi}{2}-\frac{\pi m}{n},
\]

so

\[
\lambda_a
=-e^{2i\theta}
=e^{-2\pi i m/n}.
\]

It follows that

\[
\lambda_a^n=1
\]

and

\[
\boxed{
U_{R_{n,m}}^n=I
}.
\]

This is not a numerical fit. It is an analytic consequence derived
directly from the scattering coefficients used in the current program.

\subsection{6.3 Closed Form of the System B Time
Series}\label{closed-form-of-the-system-b-time-series}

The population difference is

\[
S_j=
\frac{2\operatorname{Re}
\left(
\langle X_0,Y_0\rangle\lambda_a^j
\right)}
{\|X_0\|^2+\|Y_0\|^2}.
\]

Under the principal condition \(\phi=0\) and \(s_0=0.01\),

\[
\frac{2\langle X_0,Y_0\rangle}
{\|X_0\|^2+\|Y_0\|^2}
=0.02,
\]

and therefore

\[
\boxed{
S_j
=0.02\cos\left[
j\left(
\pi+2\arcsin\sqrt R
\right)
\right]
}.
\]

This equation shows that the iterated scattering in the original program
and the Chebyshev recurrence used for the multiprecision calculation
evaluate the same time series.

\subsection{6.4 Low Sensitivity to Initial
Values}\label{low-sensitivity-to-initial-values}

For a general initial state,

\[
S_j=C\cos(j\omega+\varphi_0),
\qquad
\omega=\arg\lambda_a.
\]

The amplitude \(C\) and initial phase \(\varphi_0\) depend on the
initial state. However,

\[
\omega(R)=\pi+2\arcsin\sqrt R
\]

is determined by the operator alone. Consequently, changing the initial
amplitude, initial phase, wave-packet shape, or harmonic content does
not move

\[
R_{n,m}.
\]

What changes is the peak height, background, and visibility.

\subsection{6.5 The Numbers 31, 124, and
248}\label{the-numbers-31-124-and-248}

At \(R_{124,23}\),

\[
\lambda_a^{31}=i,
\qquad
\lambda_a^{62}=-1,
\qquad
\lambda_a^{93}=-i,
\qquad
\lambda_a^{124}=1.
\]

Thus,

\[
31\times4=124
\]

is one complete operator period, and

\[
124\times2=248
\]

is the number of samples used by the gray metric to read two complete
periods.

The number 248 was not inserted into the scattering coefficients. It
appeared as the time window that reads the 124-cycle closed orbit twice.

\begin{center}\rule{0.5\linewidth}{0.5pt}\end{center}

\section{7. Program Structure}\label{program-structure}

\subsection{7.1 From Reflection Probability to Complex Scattering
Amplitudes}\label{from-reflection-probability-to-complex-scattering-amplitudes}

The following excerpt is the core of the minimal System B implementation
v5:

\begin{Shaded}
\begin{Highlighting}[]
\KeywordTok{def}\NormalTok{ scattering\_coefficients(reflection\_rate: }\BuiltInTok{float}\NormalTok{):}
    \ControlFlowTok{if}\NormalTok{ reflection\_rate }\OperatorTok{\textless{}} \FloatTok{0.0} \KeywordTok{or}\NormalTok{ reflection\_rate }\OperatorTok{\textgreater{}} \FloatTok{1.0}\NormalTok{:}
        \ControlFlowTok{raise} \PreprocessorTok{ValueError}\NormalTok{(}\StringTok{"reflection rate must be in [0, 1]"}\NormalTok{)}

\NormalTok{    delta\_f }\OperatorTok{=} \FloatTok{2.0} \OperatorTok{*}\NormalTok{ math.asin(math.sqrt(reflection\_rate))}
\NormalTok{    half\_delta }\OperatorTok{=} \FloatTok{0.5} \OperatorTok{*}\NormalTok{ delta\_f}
\NormalTok{    phase }\OperatorTok{=} \BuiltInTok{complex}\NormalTok{(math.cos(half\_delta), math.sin(half\_delta))}

\NormalTok{    t }\OperatorTok{=}\NormalTok{ phase }\OperatorTok{*}\NormalTok{ math.cos(half\_delta)}
\NormalTok{    r }\OperatorTok{=} \OperatorTok{{-}}\OtherTok{1j} \OperatorTok{*}\NormalTok{ phase }\OperatorTok{*}\NormalTok{ math.sin(half\_delta)}
    \ControlFlowTok{return} \BuiltInTok{complex}\NormalTok{(t), }\BuiltInTok{complex}\NormalTok{(r), }\BuiltInTok{abs}\NormalTok{(t) }\OperatorTok{**} \DecValTok{2}\NormalTok{, }\BuiltInTok{abs}\NormalTok{(r) }\OperatorTok{**} \DecValTok{2}
\end{Highlighting}
\end{Shaded}

The only input to this function is \(R\). Neither 137, 128, \(\alpha\),
\(N(R)\), nor a finite-order root enters the calculation of the
scattering amplitudes.

Source:\\
\href{../20260715/run_minimal_system_B_gray_direct_check_v5.py\#L48-L56}{run\_minimal\_system\_B\_gray\_direct\_check\_v5.py}

\subsection{7.2 Scattering Update along the Time
Direction}\label{scattering-update-along-the-time-direction}

The core operation that generates the peaks is

\begin{Shaded}
\begin{Highlighting}[]
\NormalTok{a, b }\OperatorTok{=}\NormalTok{ normalize\_pair(}
\NormalTok{    r }\OperatorTok{*}\NormalTok{ a }\OperatorTok{+}\NormalTok{ t }\OperatorTok{*}\NormalTok{ b,}
\NormalTok{    t }\OperatorTok{*}\NormalTok{ a }\OperatorTok{+}\NormalTok{ r }\OperatorTok{*}\NormalTok{ b,}
\NormalTok{)}
\end{Highlighting}
\end{Shaded}

The iterative core is

\begin{Shaded}
\begin{Highlighting}[]
\ControlFlowTok{for}\NormalTok{ step }\KeywordTok{in} \BuiltInTok{range}\NormalTok{(steps }\OperatorTok{+} \DecValTok{1}\NormalTok{):}
\NormalTok{    p\_a }\OperatorTok{=} \BuiltInTok{abs}\NormalTok{(a) }\OperatorTok{**} \DecValTok{2}
\NormalTok{    p\_b }\OperatorTok{=} \BuiltInTok{abs}\NormalTok{(b) }\OperatorTok{**} \DecValTok{2}
\NormalTok{    q }\OperatorTok{=}\NormalTok{ p\_a }\OperatorTok{+}\NormalTok{ p\_b}
\NormalTok{    s\_value }\OperatorTok{=}\NormalTok{ (p\_a }\OperatorTok{{-}}\NormalTok{ p\_b) }\OperatorTok{/}\NormalTok{ q}

    \CommentTok{\# readout of the current state}
\NormalTok{    prefix\_sums.append(prefix\_sums[}\OperatorTok{{-}}\DecValTok{1}\NormalTok{] }\OperatorTok{+}\NormalTok{ s\_value)}

    \ControlFlowTok{if}\NormalTok{ step }\OperatorTok{\textgreater{}=}\NormalTok{ steps:}
        \ControlFlowTok{break}

\NormalTok{    a, b }\OperatorTok{=}\NormalTok{ normalize\_pair(}
\NormalTok{        r }\OperatorTok{*}\NormalTok{ a }\OperatorTok{+}\NormalTok{ t }\OperatorTok{*}\NormalTok{ b,}
\NormalTok{        t }\OperatorTok{*}\NormalTok{ a }\OperatorTok{+}\NormalTok{ r }\OperatorTok{*}\NormalTok{ b,}
\NormalTok{    )}
\end{Highlighting}
\end{Shaded}

Source:\\
\href{../20260715/run_minimal_system_B_gray_direct_check_v5.py\#L81-L136}{run\_minimal\_system\_B\_gray\_direct\_check\_v5.py}

From the second iteration onward, the output \(A',B'\) of the preceding
scattering event is used as the input to the next application of the
same scattering map. In the present code, the variable \texttt{step}
carries the time-development direction.

\subsection{7.3 Gray Error}\label{gray-error}

The evaluation function is

\begin{Shaded}
\begin{Highlighting}[]
\KeywordTok{def}\NormalTok{ prefix\_error(prefix\_sums, s\_min, s\_max):}
\NormalTok{    n }\OperatorTok{=} \BuiltInTok{len}\NormalTok{(prefix\_sums) }\OperatorTok{{-}} \DecValTok{1}
\NormalTok{    total }\OperatorTok{=}\NormalTok{ prefix\_sums[n]}
\NormalTok{    s\_mean }\OperatorTok{=}\NormalTok{ total }\OperatorTok{/}\NormalTok{ n}
\NormalTok{    s\_amp }\OperatorTok{=} \FloatTok{0.5} \OperatorTok{*}\NormalTok{ (s\_max }\OperatorTok{{-}}\NormalTok{ s\_min)}

\NormalTok{    mid }\OperatorTok{=}\NormalTok{ n }\OperatorTok{//} \DecValTok{2}
\NormalTok{    first\_mean }\OperatorTok{=}\NormalTok{ prefix\_sums[mid] }\OperatorTok{/}\NormalTok{ mid}
\NormalTok{    second\_mean }\OperatorTok{=}\NormalTok{ (total }\OperatorTok{{-}}\NormalTok{ prefix\_sums[mid]) }\OperatorTok{/}\NormalTok{ (n }\OperatorTok{{-}}\NormalTok{ mid)}
\NormalTok{    s\_drift }\OperatorTok{=} \BuiltInTok{abs}\NormalTok{(second\_mean }\OperatorTok{{-}}\NormalTok{ first\_mean)}

\NormalTok{    gray\_error }\OperatorTok{=}\NormalTok{ (}
        \BuiltInTok{abs}\NormalTok{(s\_mean)}
        \OperatorTok{+} \BuiltInTok{abs}\NormalTok{(s\_amp }\OperatorTok{{-}} \FloatTok{0.02}\NormalTok{)}
        \OperatorTok{+}\NormalTok{ s\_drift}
\NormalTok{    )}
    \ControlFlowTok{return}\NormalTok{ gray\_error, s\_mean, s\_amp, s\_drift}
\end{Highlighting}
\end{Shaded}

Source:\\
\href{../20260715/run_minimal_system_B_gray_direct_check_v5.py\#L63-L78}{run\_minimal\_system\_B\_gray\_direct\_check\_v5.py}

No value 137 or 128 appears here either. The function evaluates only the
mean, amplitude, and temporal drift of a periodic orbit.

\subsection{7.4 Analytic Eigenvalues and Finite-Order
Roots}\label{analytic-eigenvalues-and-finite-order-roots}

The following analytic code was added after the numerical peaks had been
found:

\begin{Shaded}
\begin{Highlighting}[]
\KeywordTok{def}\NormalTok{ exact\_eigenvalues(reflection\_rate: }\BuiltInTok{float}\NormalTok{):}
\NormalTok{    theta }\OperatorTok{=}\NormalTok{ math.asin(math.sqrt(reflection\_rate))}
\NormalTok{    lambda\_symmetric }\OperatorTok{=} \FloatTok{1.0} \OperatorTok{+} \OtherTok{0.0j}
\NormalTok{    lambda\_antisymmetric }\OperatorTok{=} \OperatorTok{{-}}\BuiltInTok{complex}\NormalTok{(}
\NormalTok{        math.cos(}\FloatTok{2.0} \OperatorTok{*}\NormalTok{ theta),}
\NormalTok{        math.sin(}\FloatTok{2.0} \OperatorTok{*}\NormalTok{ theta),}
\NormalTok{    )}
    \ControlFlowTok{return}\NormalTok{ lambda\_symmetric, lambda\_antisymmetric}


\KeywordTok{def}\NormalTok{ resonant\_R(period\_n: }\BuiltInTok{int}\NormalTok{, mode\_m: }\BuiltInTok{int}\NormalTok{) }\OperatorTok{{-}\textgreater{}} \BuiltInTok{float}\NormalTok{:}
    \ControlFlowTok{return}\NormalTok{ math.cos(math.pi }\OperatorTok{*}\NormalTok{ mode\_m }\OperatorTok{/}\NormalTok{ period\_n) }\OperatorTok{**} \DecValTok{2}
\end{Highlighting}
\end{Shaded}

Source:\\
\href{../20260717/phase5_eigenphase_resonance_v2.py\#L74-L103}{phase5\_eigenphase\_resonance\_v2.py}

This code did not insert the roots in order to create the peaks. It was
written afterward to identify the peaks from the full-range sweep as
operator-eigenvalue roots.

\subsection{7.5 Closed-Form Multiprecision
Recurrence}\label{closed-form-multiprecision-recurrence}

Rather than directly evaluating

\[
S_j=0.02\cos(j\omega),
\]

the multiprecision calculation used a Chebyshev recurrence to control
rounding error:

\begin{Shaded}
\begin{Highlighting}[]
\NormalTok{phase }\OperatorTok{=}\NormalTok{ mp.pi }\OperatorTok{+} \DecValTok{2} \OperatorTok{*}\NormalTok{ mp.asin(mp.sqrt(reflection\_rate))}
\NormalTok{phase\_cos }\OperatorTok{=}\NormalTok{ mp.cos(phase)}

\NormalTok{cosine\_previous }\OperatorTok{=}\NormalTok{ mp.mpf(}\DecValTok{1}\NormalTok{)}
\NormalTok{cosine\_current }\OperatorTok{=}\NormalTok{ phase\_cos}

\ControlFlowTok{for}\NormalTok{ sample\_index }\KeywordTok{in} \BuiltInTok{range}\NormalTok{(sample\_count):}
    \ControlFlowTok{if}\NormalTok{ sample\_index }\OperatorTok{==} \DecValTok{0}\NormalTok{:}
\NormalTok{        cosine\_value }\OperatorTok{=}\NormalTok{ cosine\_previous}
    \ControlFlowTok{elif}\NormalTok{ sample\_index }\OperatorTok{==} \DecValTok{1}\NormalTok{:}
\NormalTok{        cosine\_value }\OperatorTok{=}\NormalTok{ cosine\_current}
    \ControlFlowTok{else}\NormalTok{:}
\NormalTok{        cosine\_previous, cosine\_current }\OperatorTok{=}\NormalTok{ (}
\NormalTok{            cosine\_current,}
            \DecValTok{2} \OperatorTok{*}\NormalTok{ phase\_cos }\OperatorTok{*}\NormalTok{ cosine\_current }\OperatorTok{{-}}\NormalTok{ cosine\_previous,}
\NormalTok{        )}
\NormalTok{        cosine\_value }\OperatorTok{=}\NormalTok{ cosine\_current}

\NormalTok{    s\_value }\OperatorTok{=}\NormalTok{ mp.mpf(}\StringTok{"0.02"}\NormalTok{) }\OperatorTok{*}\NormalTok{ cosine\_value}
\end{Highlighting}
\end{Shaded}

Source:\\
\href{run_two_physical_roots_multiprecision_v1.py\#L64-L106}{run\_two\_physical\_roots\_multiprecision\_v1.py}

This is not a separate model. It is a closed-form calculation of the
same time series derived from the original two-channel exchange update.

\subsection{7.6 Vector Wave-Packet Update in System
A}\label{vector-wave-packet-update-in-system-a}

System A applies the same scattering kernel to wave-packet vectors
rather than scalar amplitudes:

\begin{Shaded}
\begin{Highlighting}[]
\ControlFlowTok{for}\NormalTok{ collision }\KeywordTok{in} \BuiltInTok{range}\NormalTok{(max\_collision }\OperatorTok{+} \DecValTok{1}\NormalTok{):}
\NormalTok{    rows.append(}
\NormalTok{        row\_for\_state(..., }\StringTok{"A\_channel"}\NormalTok{, a, ...)}
\NormalTok{    )}
\NormalTok{    rows.append(}
\NormalTok{        row\_for\_state(..., }\StringTok{"B\_channel"}\NormalTok{, b, ...)}
\NormalTok{    )}

    \ControlFlowTok{if}\NormalTok{ collision }\OperatorTok{\textgreater{}=}\NormalTok{ max\_collision:}
        \ControlFlowTok{break}

\NormalTok{    a\_next }\OperatorTok{=}\NormalTok{ src.normalize(r }\OperatorTok{*}\NormalTok{ a }\OperatorTok{+}\NormalTok{ t }\OperatorTok{*}\NormalTok{ b)}
\NormalTok{    b\_next }\OperatorTok{=}\NormalTok{ src.normalize(t }\OperatorTok{*}\NormalTok{ a }\OperatorTok{+}\NormalTok{ r }\OperatorTok{*}\NormalTok{ b)}
\NormalTok{    a, b }\OperatorTok{=}\NormalTok{ a\_next, b\_next}
\end{Highlighting}
\end{Shaded}

Source:\\
\href{../20260715/run_system_A_localization_exchange_R_sweep_preliminary_v1.py\#L718-L737}{run\_system\_A\_localization\_exchange\_R\_sweep\_preliminary\_v1.py}

Systems A and B use different observables, but their time-development
kernel \(U_R\) is identical.

\subsection{7.7 Collision Trigger and the Scope of Iterated
Scattering}\label{collision-trigger-and-the-scope-of-iterated-scattering}

The complete interaction separates into two stages:

\begin{enumerate}
\def\labelenumi{\arabic{enumi}.}
\tightlist
\item
  The A and B wave packets approach within one cell, triggering the
  first scattering event.
\item
  The post-scattering outputs \(A',B'\) are updated into the next cycle
  of iterated scattering.
\end{enumerate}

The current System A and System B programs directly implement the second
stage. They do not implement the trajectories by which wave packets
approach through space until the first-collision condition is satisfied.

The result of this paper is therefore not a completed model of the
entire interaction process. It is the identification that \textbf{the
iterated exchange-scattering kernel after collision onset is sufficient
to generate finite-order resonance}.

\begin{center}\rule{0.5\linewidth}{0.5pt}\end{center}

\section{8. Audit for Hidden
Parameters}\label{audit-for-hidden-parameters}

\subsection{8.1 Classification of the Code
Series}\label{classification-of-the-code-series}

The existing code contains three distinct classes with different
purposes.

{\def\LTcaptype{none} % do not increment counter
\begin{longtable}[]{@{}
  >{\raggedright\arraybackslash}p{(\linewidth - 4\tabcolsep) * \real{0.3333}}
  >{\raggedright\arraybackslash}p{(\linewidth - 4\tabcolsep) * \real{0.3333}}
  >{\raggedright\arraybackslash}p{(\linewidth - 4\tabcolsep) * \real{0.3333}}@{}}
\toprule\noalign{}
\begin{minipage}[b]{\linewidth}\raggedright
Code class
\end{minipage} & \begin{minipage}[b]{\linewidth}\raggedright
Known 137/128 values
\end{minipage} & \begin{minipage}[b]{\linewidth}\raggedright
Role
\end{minipage} \\
\midrule\noalign{}
\endhead
\bottomrule\noalign{}
\endlastfoot
Early exploration and comparison code & Included & Compare robustness to
harmonics and initial conditions at candidate points \\
Uniform System B v5 sweep & Not included in sweep evaluation & Discover
candidate peaks from an equally spaced grid \\
Analytic-root and multiprecision code & Roots explicit & Verify
agreement between discovered peaks and analytic roots \\
\end{longtable}
}

The early System A exploration code and early System B code explicitly
contained the 137- and 128-corresponding points as fixed probes. Those
calculations alone are therefore not used as evidence of independent
discovery.

In contrast, the System B v5 sweep mode generated a Decimal equally
spaced sequence from command-line minimum, maximum, and spacing
parameters, and applied the same evaluation function at every point:

\begin{Shaded}
\begin{Highlighting}[]
\NormalTok{r\_values }\OperatorTok{=}\NormalTok{ decimal\_range(}
\NormalTok{    args.min\_R,}
\NormalTok{    args.max\_R,}
\NormalTok{    args.delta\_R,}
\NormalTok{)}

\ControlFlowTok{for}\NormalTok{ r\_text }\KeywordTok{in}\NormalTok{ r\_values:}
\NormalTok{    payload }\OperatorTok{=}\NormalTok{ probe\_r(}
        \BuiltInTok{float}\NormalTok{(r\_text),}
\NormalTok{        r\_text,}
\NormalTok{        args.steps,}
\NormalTok{        args.phi\_mode,}
\NormalTok{        args.min\_steps,}
\NormalTok{        args.early\_stop\_patience,}
\NormalTok{    )}
\end{Highlighting}
\end{Shaded}

Source:\\
\href{../20260715/run_minimal_system_B_gray_direct_check_v5.py\#L228-L270}{run\_minimal\_system\_B\_gray\_direct\_check\_v5.py}

A default \(R\) for single-point execution remains in the source, but it
is not referenced when sweep mode is invoked with \texttt{min-R},
\texttt{max-R}, and \texttt{delta-R}.

\subsection{8.2 Chronology of Discovery and
Verification}\label{chronology-of-discovery-and-verification}

The evidential order is:

\begin{enumerate}
\def\labelenumi{\arabic{enumi}.}
\tightlist
\item
  Observe localization transfer and robustness to harmonics in System A.
\item
  Construct the metastable interface in System B.
\item
  Perform the full \(R\) sweep with the explicit aim of testing the
  \(\alpha\) correspondence.
\item
  Discover the principal peaks on a \(10^{-7}\) grid.
\item
  Perform root-centered precision sweeps.
\item
  Find that the peak centers coincide with \(R_{124,23}\) and
  \(R_{122,23}\).
\item
  Derive the general roots \(R_{n,m}\) from the operator eigenvalue.
\item
  Identify \(R_{620,117}\) near the high-energy observed value.
\item
  Confirm precision-dependent growth of the depth with 50- and 80-digit
  calculations.
\item
  Only then search the literature and recognize connections to exact
  Floquet recurrence, Farey structure, and \(E_8\).
\end{enumerate}

This chronology establishes that the peaks were not generated by
embedding a previously known cyclotomic root or the \(31\times8\)
structure of \(E_8\) in the numerical code.

\subsection{8.3 Reproduction Conditions}\label{reproduction-conditions}

The conditions of the full-range sweep are preserved in the CSV file

\href{../20260715/minimal_system_B_gray_bugcheck_result_v1/direct_depth_probe_v5_sweep_control/high_to_ext_full_delta1e-7_stats_v5.csv}{high\_to\_ext\_full\_delta1e-7\_stats\_v5.csv}.

The principal parameters are:

{\def\LTcaptype{none} % do not increment counter
\begin{longtable}[]{@{}lr@{}}
\toprule\noalign{}
Item & Value \\
\midrule\noalign{}
\endhead
\bottomrule\noalign{}
\endlastfoot
\(R_{\min}\) & 0.68660290255614798 \\
\(R_{\max}\) & 0.70246536756843059 \\
\(\Delta R\) & \(10^{-7}\) \\
Number of points & 158,625 \\
steps & 1024 \\
min steps & 256 \\
patience & 20 \\
phi & 0 \\
\end{longtable}
}

The inputs, precision, number of points, raw JSON data, and figures for
the local and multiprecision sweeps are likewise preserved in their
result directories.

\begin{center}\rule{0.5\linewidth}{0.5pt}\end{center}

\section{9. The Two Physical Models and Resonance
Visibility}\label{the-two-physical-models-and-resonance-visibility}

\subsection{9.1 Harmonics Do Not Create the Roots; They Make the Roots
Visible}\label{harmonics-do-not-create-the-roots-they-make-the-roots-visible}

Let the internal degrees of freedom of a wave packet be \(\mathcal H\).
Where the same exchange law acts on every internal component,

\[
U_R^{\mathrm{total}}
=U_R\otimes I_{\mathcal H}.
\]

The harmonic type, amplitude, phase, and wavelength change the initial
vector in \(\mathcal H\), but they do not change the two-channel
eigenvalues

\[
1,\qquad -e^{2i\theta}.
\]

Additional harmonics are therefore not the cause that creates the
finite-order roots. They are degrees of freedom that couple the
antisymmetric eigensector to localization or population-difference
observables and thereby make the resonance visible.

The exact conditions under which the resonance does not appear in the
System B population difference are

\[
Y_0=0
\]

or

\[
\langle X_0,Y_0\rangle=0.
\]

In these cases, the antisymmetric eigenphase exists in the operator but
is absent from the selected readout.

\subsection{9.2 Why Different Models Read the Same
Region}\label{why-different-models-read-the-same-region}

System A reads localization, harmonic order, and transfer ratio through

\[
\mathcal O_A(U_R^j\Psi_0).
\]

System B reads the A/B population difference, gray error, and metastable
depth through

\[
\mathcal O_B(U_R^j\Psi_0).
\]

The observation functions differ, but the time-development operator is
the same. The different physical models respond in the same \(R\) region
because their common exchange-scattering kernel has the same eigenphase.

This is the point of contact with physics in the present study. We did
not begin by enumerating abstract cyclotomic roots. Experimental code
designed to investigate wave-packet localization transfer and metastable
selection detected the same operator resonance.

\begin{center}\rule{0.5\linewidth}{0.5pt}\end{center}

\section{10. Connections to Prior Work Recognized after the
Discovery}\label{connections-to-prior-work-recognized-after-the-discovery}

\subsection{10.1 Chronology of the Literature
Search}\label{chronology-of-the-literature-search}

The works discussed below did not guide the full-range sweep or the
program design. We searched the literature only after identifying
\(R_{124,23}\), \(R_{122,23}\), and \(R_{620,117}\) through numerical
calculation and eigenvalue analysis.

This chronology is stated explicitly for two reasons:

\begin{enumerate}
\def\labelenumi{\arabic{enumi}.}
\tightlist
\item
  to show that the numerical experiment was not a reproduction obtained
  by embedding known cyclotomic roots or an \(E_8\) decomposition in the
  code; and
\item
  to distinguish what is known as a general phenomenon from the
  mathematical model and numerical discovery specific to this study.
\end{enumerate}

\subsection{10.2 Exact State-Independent Quantum
Recurrence}\label{exact-state-independent-quantum-recurrence}

Anand et al.~classified exact state-independent recurrence times from
the cyclotomic spectrum of finite-dimensional Floquet unitary operators
{[}5{]}. The present relation

\[
U_R^n=I
\]

is a concrete two-channel instance connected to this general class of
exact recurrence.

We therefore do not claim to have discovered the general concept of a
finite-order unitary operator. The contribution is that a wave-packet
numerical experiment aimed at the fine-structure constant independently
reached this structure and detected it in both localization-transfer and
metastable-interface models.

\subsection{10.3 Farey Sequences and Resonance
Organization}\label{farey-sequences-and-resonance-organization}

Tomás described the organization of physical resonance diagrams by Farey
sequences of rational numbers {[}6{]}. The present roots are classified
by

\[
\frac{m}{n}.
\]

Farey adjacency can therefore organize the closely spaced high-order
roots on the high-energy side. That work, however, does not treat the
exchange-scattering map, the fine-structure-constant correspondence,
localization transfer, or the cat model used here.

\subsection{\texorpdfstring{10.4 The \(31\times8\) Structure of
\(E_8\)}{10.4 The 31\textbackslash times8 Structure of E\_8}}\label{the-31times8-structure-of-e_8}

After obtaining

\[
124=31\times4,
\qquad
248=31\times4\times2,
\]

we investigated a possible connection to the 248 dimensions of \(E_8\)
and its decomposition into 31 orthogonal eight-dimensional subspaces.

Thompson and later work treat structures associated with a decomposition
of \(E_8\) into 31 orthogonal eight-dimensional Cartan subalgebras
{[}7,8{]}. The \(E_8\) root system also separates into 112 integer-type
roots and 128 half-integer-type roots {[}9{]}.

These connections were noticed only after the finite-order-root
discovery of July 18, 2026. The numbers 31, 128, and 248 were not
selected in advance to tune the iteration count or exchange coefficient.

At present, no explicit map has been constructed from

\[
31\text{ time sectors}\times
4\text{ phases}\times
2\text{ channels}
\]

to the 31 eight-dimensional subspaces of \(E_8\). The \(E_8\) connection
is therefore retained as a structural hypothesis.

\subsection{10.5 Exceptional Algebras and the Fine-Structure
Constant}\label{exceptional-algebras-and-the-fine-structure-constant}

El Naschie proposed a derivation connecting the \(8\times31=248\)
structure of \(E_8\) to the fine-structure constant {[}10{]}. Singh
discussed \(1/137\) using octonionic spacetime and an exceptional Jordan
algebra {[}11{]}.

Neither line of work was the starting point of this study. We
deliberately searched for the fine-structure constant, but did not use
\(E_8\) or an exceptional algebra to design the numerical peaks. Their
relation to the present result was investigated only after obtaining the
finite-order roots and the \(31\times4\times2\) structure.

\begin{center}\rule{0.5\linewidth}{0.5pt}\end{center}

\section{11. Discussion}\label{discussion}

\subsection{11.1 What Did We Discover?}\label{what-did-we-discover}

We initially expected that a closed exchange system could directly
explain the fine-structure constant.

The numerical experiments did return sharp peaks near 137 and 128. Their
direct cause, however, was not the assumed electromagnetic-coupling
correspondence

\[
G_R=e
\]

itself. The direct cause is the finite-order exchange resonance

\[
\lambda_a^n=1.
\]

Thus, this study has not completed a direct derivation of \(\alpha\). In
the course of searching for the fine-structure constant, it discovered
that the iterated exchange-scattering program detects exact finite-order
recurrence.

Because general exact recurrence is present in prior work, the general
phenomenon is a rediscovery of a different phenomenon. The following
points are nevertheless specific to this study:

\begin{itemize}
\tightlist
\item
  it began from a wave-packet localization-transfer model;
\item
  it also began from a weak-readout/strong-selection model of two
  metastable states;
\item
  both models share the same exchange-scattering kernel;
\item
  multiple even-order roots appear near the physical readout regions 137
  and 128; and
\item
  source code, raw data, full-range figures, local figures, and
  multiprecision results are preserved in reproducible form.
\end{itemize}

\subsection{\texorpdfstring{11.2 Why Are the Roots Near
\(\alpha\)?}{11.2 Why Are the Roots Near \textbackslash alpha?}}\label{why-are-the-roots-near-alpha}

This is the most important question remaining after the result of this
paper.

On the low-energy side,

\[
N(R_{124,23})
=137.036042914598545\ldots
\]

is close to

\[
\alpha^{-1}(0)
=137.035999177(21).
\]

On the high-energy side,

\[
N(R_{620,117})
=128.947864735670559\ldots
\]

is close to

\[
\alpha^{-1}(M_Z^2)
=128.946\pm0.015.
\]

There are countably infinitely many finite-order roots

\[
R_{n,m}=\cos^2\left(\frac{\pi m}{n}\right).
\]

Their arithmetic density can place high-order roots near an arbitrary
real value and may account for part of the proximity.

However, this study did not merely select arbitrary roots after the
fact:

\begin{itemize}
\tightlist
\item
  the original full-range sweep produced \(R_{124,23}\) and
  \(R_{122,23}\) as its principal peaks;
\item
  the low-energy root has the short \(31\times4\) temporal structure;
\item
  an even root \(R_{620,117}\) with infinite depth under the present
  gray metric lies near the high-energy observed value; and
\item
  different wave-packet models responded to the same exchange-scattering
  region.
\end{itemize}

The proximity therefore cannot yet be declared either a mere coincidence
or a physical necessity.

\subsection{11.3 The Central Open
Problem}\label{the-central-open-problem}

The central open problem of this paper is stated without ambiguity:

\begin{quote}
\textbf{We began this search expecting that a closed exchange system
could explain the fine-structure constant \(\alpha\). Within what has
been established by the present analysis, however, we did not derive
\(\alpha\) itself; we merely rediscovered a different finite-order
resonance phenomenon connected to the general theory of exact
recurrence. Yet the extremely sharp peaks of that finite-order resonance
actually occur near the two physical low- and high-energy values of
\(\alpha^{-1}\). Why this proximity occurs remains unexplained. Is it an
accidental consequence of arithmetic density, or is it another necessity
imposed by four-dimensional discrete geometry, \(E_8\)-type symmetry, a
physical selection rule of exchange scattering, or an
as-yet-unidentified structure? Distinguishing these alternatives is the
central task of future work.}
\end{quote}

The required discriminating calculations are:

\begin{enumerate}
\def\labelenumi{\arabic{enumi}.}
\tightlist
\item
  enumerate all roots below a fixed maximum order \(n_{\max}\);
\item
  calculate how often an even-order root at least as close to the
  observed value occurs near randomized reference values;
\item
  test whether localization, entropy, phase correlation, and transfer
  observables select the same roots independently of the gray metric;
\item
  randomize the initial state, harmonics, and time window to separate
  root location from visibility;
\item
  construct an explicit map connecting the four-dimensional 137-cell
  count, the 128 outer sector, the \(31\times8\) decomposition, and the
  exchange operator; and
\item
  implement the same \(U_R\) in a physical optical coupler,
  discrete-time quantum walk, or two-level Floquet system.
\end{enumerate}

\begin{center}\rule{0.5\linewidth}{0.5pt}\end{center}

\section{12. Conclusion}\label{conclusion}

This study identified the direct cause of the sharp peaks observed near
the fine-structure-constant inverse values 137 and 128.

The peaks form where the antisymmetric eigenvalue

\[
\lambda_a=-e^{2i\arcsin\sqrt R}
\]

of

\[
U_R=
\begin{pmatrix}
r&t\\t&r
\end{pmatrix}
\]

has finite order, namely at

\[
R_{n,m}
=\cos^2\left(\frac{\pi m}{n}\right).
\]

The principal low-energy peak is \(R_{124,23}\), the former high-energy
principal candidate is \(R_{122,23}\), and \(R_{620,117}\) lies near the
high-energy observed value.

With 50- and 80-digit arithmetic, the depth increases with numerical
precision. In ideal arithmetic, the gray error vanishes at an even
fundamental-order root and the resonance depth becomes infinite.

The result is not the first presentation of general quantum recurrence.
It is the discovery that a wave-packet numerical experiment deliberately
aimed at the fine-structure constant detects finite-order roots without
inserting the known values into its scattering kernel, and that the
phenomenon can be verified in reproducible localization-transfer and
metastable two-state models.

The fine-structure constant itself remains unexplained. A new question
has been isolated: why do extremely sharp finite-order roots occur near
the two physical values of \(\alpha^{-1}\)?

The next stage is to distinguish

\[
\text{coincidence}
\quad\text{from}\quad
\text{another necessity}.
\]

\begin{center}\rule{0.5\linewidth}{0.5pt}\end{center}

\section{References}\label{references}

\subsection{Self-References}\label{self-references}

\begin{enumerate}
\def\labelenumi{\arabic{enumi}.}
\tightlist
\item
  Noriaki Kihara,
  \href{../20260713/交換干渉散乱行列フェルミオン的衝突における低局在性・倍音移乗読出し実験仕様書\%20v1.md}{Experimental
  Specification v1 for Low-Localization and Harmonic-Transfer Readout in
  Exchange-Interference Scattering-Matrix Fermion-Like Collisions},
  2026.
\item
  Noriaki Kihara,
  \href{../20260713/exchange_scattering_matrix_fermionic_localization_exchange_preliminary_summary_en.md}{Preliminary
  Summary of Acceleration-Basis and Localization Exchange in
  Exchange-Interference Scattering-Matrix Fermion-Like Collisions},
  Zenodo Concept DOI: 10.5281/zenodo.21333766, 2026.
\item
  Noriaki Kihara,
  \href{../20260714/gray_cat_metastable_observation_preliminary_summary_en.md}{Preliminary
  Numerical Study of Gray-Cat Metastable Interfaces, Weak C Readout, and
  Strong D Observation Selection in a Closed Phase System}, Zenodo
  Concept DOI: 10.5281/zenodo.21353208, 2026.
\item
  Noriaki Kihara,
  \href{../20260715/exchange_weight_alpha_correspondence_numerical_experiment_en.md}{Selection
  of the State-Exchange Weight \(G_R=1-R\) and Candidate Correspondence
  with the Fine-Structure Constant: A Numerical Experiment}, Zenodo
  Concept DOI: 10.5281/zenodo.21396760, 2026.
\end{enumerate}

\subsection{External References}\label{external-references}

\begin{enumerate}
\def\labelenumi{\arabic{enumi}.}
\setcounter{enumi}{4}
\tightlist
\item
  Amit Anand, Dinesh Valluri, Jack Davis, and Shohini Ghose, ``Quantum
  recurrences and the arithmetic of Floquet dynamics,'' \emph{Quantum}
  \textbf{10}, 2074 (2026).
  \url{https://doi.org/10.22331/q-2026-04-20-2074}
\item
  R. Tomás, ``From Farey sequences to resonance diagrams,''
  \emph{Physical Review Special Topics---Accelerators and Beams}
  \textbf{17}, 014001 (2014).
  \url{https://doi.org/10.1103/PhysRevSTAB.17.014001}
\item
  J. G. Thompson, ``A conjugacy theorem for \(E_8\),'' \emph{Journal of
  Algebra} \textbf{38}, 525--530 (1976).
  \url{https://doi.org/10.1016/0021-8693(76)90235-0}
\item
  Alexander A. Ivanov, ``Geometry of the Thompson group,'' \emph{Journal
  of Algebra} \textbf{661}, 778--806 (2025).
  \url{https://doi.org/10.1016/j.jalgebra.2024.08.004}
\item
  Rosa Winter and Ronald van Luijk, ``The Action of the Weyl Group on
  the \(E_8\) Root System,'' \emph{Graphs and Combinatorics}
  \textbf{37}, 1965--2064 (2021).
  \url{https://doi.org/10.1007/s00373-021-02315-8}
\item
  M. S. El Naschie, ``A derivation of the fine structure constant from
  the exceptional Lie group hierarchy of the micro cosmos,''
  \emph{Chaos, Solitons \& Fractals} \textbf{36}, 819--822 (2008).
  \url{https://doi.org/10.1016/j.chaos.2007.09.020}
\item
  Tejinder P. Singh, ``Quantum theory without classical time: octonions,
  and a theoretical derivation of the fine structure constant 1/137,''
  \emph{International Journal of Modern Physics D} \textbf{30}, 2142010
  (2021). \url{https://doi.org/10.1142/S0218271821420104}
\item
  Pierre-Philippe Dechant, ``The birth of \(E_8\) out of the spinors of
  the icosahedron,'' \emph{Proceedings of the Royal Society A}
  \textbf{472}, 20150504 (2016).
  \url{https://doi.org/10.1098/rspa.2015.0504}
\item
  CODATA 2022, ``The 2022 CODATA Recommended Values of the Fundamental
  Physical Constants,'' \(\alpha^{-1}(0)=137.035999177(21)\).
  \url{https://physics.nist.gov/cuu/pdf/wallet_2022.pdf}
\item
  Alexander Keshavarzi, Daisuke Nomura, and Thomas Teubner, ``The muon
  \(g-2\) and \(\alpha(M_Z^2)\): a new data-based analysis,''
  \emph{Physical Review D} \textbf{101}, 014029 (2020).
  \url{https://doi.org/10.1103/PhysRevD.101.014029}
\end{enumerate}

\begin{center}\rule{0.5\linewidth}{0.5pt}\end{center}

\section{Appendix A. Principal Reproduction
Files}\label{appendix-a.-principal-reproduction-files}

{\def\LTcaptype{none} % do not increment counter
\begin{longtable}[]{@{}
  >{\raggedright\arraybackslash}p{(\linewidth - 2\tabcolsep) * \real{0.5000}}
  >{\raggedright\arraybackslash}p{(\linewidth - 2\tabcolsep) * \real{0.5000}}@{}}
\toprule\noalign{}
\begin{minipage}[b]{\linewidth}\raggedright
Content
\end{minipage} & \begin{minipage}[b]{\linewidth}\raggedright
File
\end{minipage} \\
\midrule\noalign{}
\endhead
\bottomrule\noalign{}
\endlastfoot
July 15, 2026 complete paper &
\href{../20260715/exchange_weight_alpha_correspondence_numerical_experiment_en.md}{Previous
paper, English} \\
Full-range sweep statistics &
\href{../20260715/minimal_system_B_gray_bugcheck_result_v1/direct_depth_probe_v5_sweep_control/high_to_ext_full_delta1e-7_stats_v5.csv}{Full-range
sweep CSV} \\
Full-range sweep figure &
\href{../20260715/system_B_full_R_sweep_full_range_depth_v1.png}{Figure
1 PNG} \\
\(R_{124,23}\) local sweep &
\href{root_centered_delta1e-12_v1/root_centered_delta1e-12_result_v1.md}{Local-sweep
result} \\
\(R_{122,23}\) resolution comparison &
\href{R122_23_root_centered_resolution_comparison_v1/R122_23_root_centered_resolution_comparison_v1.md}{Resolution-comparison
result} \\
Higher-order-root comparison &
\href{high_order_roots_centered_delta1e-12_v1/high_order_roots_centered_delta1e-12_result_v1.md}{Higher-order-root
result} \\
Multiprecision result &
\href{two_physical_roots_multiprecision_v1/two_physical_roots_multiprecision_result_v1.md}{Multiprecision
result} \\
Precision-matched ultralocal sweep &
\href{two_physical_roots_precision_matched_narrow_v1/precision_matched_narrow_result_v1.md}{Ultralocal-sweep
result} \\
Minimal System B scattering kernel &
\href{../20260715/run_minimal_system_B_gray_direct_check_v5.py}{System B
source} \\
Eigenvalue and analytic-root code &
\href{../20260717/phase5_eigenphase_resonance_v2.py}{Eigenvalue-analysis
source} \\
Multiprecision code &
\href{run_two_physical_roots_multiprecision_v1.py}{Multiprecision
source} \\
\end{longtable}
}

\begin{center}\rule{0.5\linewidth}{0.5pt}\end{center}

\section{Appendix B. Classification of
Claims}\label{appendix-b.-classification-of-claims}

{\def\LTcaptype{none} % do not increment counter
\begin{longtable}[]{@{}
  >{\raggedright\arraybackslash}p{(\linewidth - 2\tabcolsep) * \real{0.5000}}
  >{\raggedright\arraybackslash}p{(\linewidth - 2\tabcolsep) * \real{0.5000}}@{}}
\toprule\noalign{}
\begin{minipage}[b]{\linewidth}\raggedright
Claim
\end{minipage} & \begin{minipage}[b]{\linewidth}\raggedright
Classification
\end{minipage} \\
\midrule\noalign{}
\endhead
\bottomrule\noalign{}
\endlastfoot
Unitarity of \(U_R\) & Derived consequence \\
\(\lambda_s=1,\lambda_a=-e^{2i\theta}\) & Derived consequence \\
\(R_{n,m}=\cos^2(\pi m/n)\) & Derived consequence \\
\(U_{R_{n,m}}^n=I\) & Derived consequence \\
Infinite ideal gray depth at even fundamental-order roots & Derived
consequence of the present readout \\
Agreement of the full-range peaks with \(R_{124,23}\) and \(R_{122,23}\)
& Numerical fact \\
Proximity of \(R_{620,117}\) to the high-energy observed value &
Numerical fact \\
Harmonics control visibility rather than root position & Theoretical
insight from the operator decomposition \\
Identity with the \(31\times8\) decomposition of \(E_8\) & Structural
hypothesis \\
Necessity of the proximity between finite-order roots and physical
\(\alpha\) & Not derived; central open problem \\
\end{longtable}
}

\end{document}
